\section{Related Literature}
\label{sec:related}

\paragraph{Residual-based deep learning in economics.} Minimizing the residuals of equilibrium conditions has a long history in economics through projection and collocation methods \parencite{Judd1992projection, JuddMailar2011}, which represent the policy with polynomials. DEQNs \parencite{Azinovic2022DEQN} and related deep-learning methods \parencite{Maliar2021, Kahou2021} replace the polynomials with neural networks and train them on simulated states, which mitigates the curse of dimensionality. \textcite{Folini2024ClimateInCE} and \textcite{friedl2023_deepUQ} apply DEQNs to climate-economy models. Outside economics, the same idea underlies physics-informed neural networks \parencite{Raissi2019PINN}, which minimize the residuals of differential equations. A strength of residual-based methods is that they apply to equilibria that are not the solution of any single optimization problem, such as decentralized economies with market clearing.

\paragraph{Reinforcement learning and optimal control.} Dynamic programming and RL share the Bellman equation as their foundation \parencite{Bertsekas2019RLOC, Sutton2018intro_to_rl}. Deterministic policy gradients \parencite{Silver2014DPG} and their deep variants \parencite{Lillicrap2016DDPG, Fujimoto2018TD3} train a deterministic actor on the gradient of a learned critic. When the dynamics are known and differentiable, the critic of the next state can be evaluated through the model, as in stochastic value gradients \parencite{Heess2015SVG}. The delayed actor updates and target networks used here follow \textcite{Fujimoto2018TD3} and \textcite{Lillicrap2016DDPG}. RL has also been applied to climate-economy models \parencite{strnard2019deepRLworld-earth, Wolf2023rl-to-iam}, but, to my knowledge, not evaluated against a closed-form solution of a quantitative IAM.

\paragraph{Value-gradient learning.} Dual heuristic programming \parencite{Werbos1992ADP} and value-gradient learning \parencite{Fairbank2012VGL} train a critic on the gradient of the value function rather than, or in addition to, its level. Sobolev training \parencite{Czarnecki2017Sobolev} shows more generally that fitting derivatives improves the accuracy and data efficiency of neural approximations. In economics, the gradient of the value function is the vector of costates or shadow prices, which makes these methods a natural fit for problems in which policy depends on shadow prices, such as the SCC.
